\section{Experiment}

\subsection{Experiment Setup}
\subsubsection{Dataset}
We evaluate our model in two reciprocal recommendation scenarios encompassing five real-world datasets.
\begin{enumerate}
\item \textbf{Online recruitment.} We conduct the evaluation using three authentic datasets sourced from a popular Chinese online recruitment platform. These datasets are created from over 100 days of real online logs, with all user private information removed. 
The employer-employee pairs that achieved an interview are gathered as instances of positive interactions.
\item \textbf{Question answering.}
We use two public datasets from distinct sections of the popular online Q\&A website ``Stack Exchange''~\cite{paranjape2017motifs}, intended for mutual matching between questioners and answerers.
For the \textit{answerer} side, we recommend questions to users according to their historical answering records. For the \textit{questioner} side, we follow an expert finding~\cite{lin2017survey} paradigm to recommend expert users for questioners.
\end{enumerate}

Regarding data processing, we initially remove inactive users through 5-core filtering and then create a behavior sequence for each user based on chronological order. For each interaction record, we only use the sequence truncated to the current interaction time to prevent data leakage during training.
On the three recruitment datasets, we adopt the last two weeks of data as the \emph{validation} and \emph{test} sets, and the remaining data is used for \emph{training}. Besides, the datasets for questioner-answerer matching are divided based on the time span in an 8:1:1 ratio.
The overall statistics are shown in Table~\ref{tab:data_statistics}.
Our code is available at \textcolor{blue}{\url{https://github.com/RUCAIBox/ReSeq/}}.

\begin{table}[]
\caption{Statistics of the experimental datasets.}
\label{tab:data_statistics}
\resizebox{0.45\textwidth}{!}{
\setlength{\tabcolsep}{1.0mm}{
\begin{tabular}{@{}ccrrrr@{}}
\toprule
\multirow{4}{*}{\textbf{Recruitment}}
 & \textbf{Datasets} & \textbf{\#Candidates} & \textbf{\#Recruiters}  & \textbf{\#Interactions} & \textbf{Sparsity} \\ \cmidrule(l){2-6}
& Design & 12,274 & 9,141  & 139,355        & 99.88\%  \\
                             & Sale & 15,831 & 12,757 & 112,340        & 99.94\%  \\
                             & Technology & 56,620 & 48,071 & 808,376 & 99.97\%  \\ \midrule
\multirow{3}{*}{\textbf{Q\&A}}        
& \textbf{Datasets} & \textbf{\#Questioners} & \textbf{\#Answerers} & \textbf{\#Interactions} & \textbf{Sparsity} \\ \cmidrule(l){2-6}
& StackOverflow & 42,381 & 25,004  & 360,063                                     & 99.97\%  \\
                             & AskUbuntu & 6,030 & 3,415 & 69,865        & 99.66\% \\ \bottomrule
\end{tabular}}}
\end{table}

\begin{table*}[]
\caption{Performance comparison of different methods on the three datasets in the recruitment scenario. The best-performed and the second-best-performed methods are indicated in bold and underlined font, respectively. ``*'' denotes that the improvements are significant at the level of 0.01 with paired $t$-test.}
\label{tab:recruit_comparison}
\resizebox{0.82\linewidth}{!}{
\setlength{\tabcolsep}{2.7mm}{
\begin{tabular}{@{}cccccccc@{}}
\toprule
\multirow{2}{*}{Dataset}     & Perspective & \multicolumn{3}{c}{Candidates}                         & \multicolumn{3}{c}{Recruiters}                      \\ \cmidrule(l){2-2} \cmidrule(l){3-5} \cmidrule(l){6-8}
                             & Metric      & HR@5            & MRR@5           & NDCG@5          & HR@5            & MRR@5           & NDCG@5          \\ \midrule
\multirow{13}{*}{Design}     & BPR         & 0.1437          & 0.0866          & 0.1008          & 0.1300            & 0.0693          & 0.0842          \\
                             & LFRR        & 0.1465          & 0.0855          & 0.1006          & 0.1338          & 0.0694          & 0.0853          \\
                             & NeuMF       & 0.1611          & 0.1016          & 0.1163          & 0.1452          & 0.0804          & 0.0964          \\
                             & LightGCN    & 0.1496          & 0.0913          & 0.1057          & 0.1349          & 0.0732          & 0.0884          \\ \cmidrule(l){2-8}
                             & SASRec      & 0.2033          & 0.1184          & 0.1394          & 0.2529          & 0.1389          & 0.1670           \\
                             & SSE-PT      & 0.1961 & 0.1097	& 0.1311	& 0.2466	& 0.1327	& 0.1608                     \\
                             & BERT4Rec    & 0.1535          & 0.0874          & 0.1037          & 0.2027          & 0.1044          & 0.1286          \\
                             & FMLP-Rec     & 0.1940           & 0.1194          & 0.1379          & \underline{0.2845}    & \underline{0.1571}    & \underline{0.1886}    \\ \cmidrule(l){2-8}
                             & PJFNN       & 0.1056          & 0.0621          & 0.0728          & 0.0604          & 0.0272          & 0.0353          \\
                             & BPJFNN      & 0.1755          & 0.1010           & 0.1194          & 0.1589          & 0.0766          & 0.0968          \\
                             & IPJF        & 0.1778          & 0.0943          & 0.1149          & 0.1661          & 0.0835          & 0.1038          \\
                             & PJFFF       & 0.1839          & 0.1052          & 0.1247          & 0.1235          & 0.0595          & 0.0752          \\
                             & DPGNN       & \underline{0.2422}    & \underline{0.1396}    & \underline{0.1649}    & 0.2317          & 0.1273          & 0.1530           \\ \cmidrule(l){2-8}
                             & ReSeq        & \textbf{0.4435*} & \textbf{0.2554*} & \textbf{0.3020*}  & \textbf{0.3722*} & \textbf{0.1994*} & \textbf{0.2420*}  \\ \midrule
\multirow{13}{*}{Sale}       & BPR         & 0.1071          & 0.0613          & 0.0726          & 0.0929          & 0.0513          & 0.0615          \\
                             & LFRR        & 0.1262          & 0.0730           & 0.0862          & 0.1083          & 0.0602          & 0.0721          \\
                             & NeuMF       & 0.1367          & 0.0828          & 0.0962          & 0.1113          & 0.0639          & 0.0756          \\
                             & LightGCN    & 0.1447          & 0.0832          & 0.0984          & 0.1228          & 0.0691          & 0.0823          \\ \cmidrule(l){2-8}
                             & SASRec      & 0.1707          & 0.0932          & 0.1123          & 0.1494          & 0.0807          & 0.0976          \\
                             & SSE-PT      & 0.1592	& 0.0912	& 0.1080	& 0.1276	& 0.0648	& 0.0802 \\
                             & BERT4Rec    & 0.1303          & 0.0685          & 0.0837          & 0.1143          & 0.0579          & 0.0718          \\
                             & FMLP-Rec     & 0.1758          & 0.1001          & 0.1188          & 0.1384          & 0.0742          & 0.0901          \\ \cmidrule(l){2-8}
                             & PJFNN       & 0.1097          & 0.0650           & 0.0760           & 0.0568          & 0.0269          & 0.0342          \\
                             & BPJFNN      & 0.1355          & 0.0767          & 0.0912          & 0.1130           & 0.0549          & 0.0692          \\
                             & IPJF        & 0.1092          & 0.0541          & 0.0677          & 0.1177          & 0.0629          & 0.0763          \\
                             & PJFFF       & 0.1564          & 0.0876          & 0.1046          & 0.1052          & 0.0524          & 0.0654          \\
                             & DPGNN       & \underline{0.1871}    & \underline{0.1026}    & \underline{0.1234}    & \underline{0.1703}    & \underline{0.0930}     & \underline{0.1121}    \\ \cmidrule(l){2-8}
                             & ReSeq        & \textbf{0.2712*} & \textbf{0.1499*} & \textbf{0.1798*} & \textbf{0.2434*} & \textbf{0.1335*} & \textbf{0.1606*} \\ \midrule
\multirow{13}{*}{Technology} & BPR         & 0.2733          & 0.1778          & 0.2015          & 0.2659          & 0.1627          & 0.1883          \\
                             & LFRR        & 0.2270           & 0.1420           & 0.1631          & 0.2151          & 0.1232          & 0.1459          \\
                             & NeuMF       & 0.2796          & 0.1821          & 0.2063          & 0.2711          & 0.1662          & 0.1922          \\
                             & LightGCN    & 0.2775          & 0.1823          & 0.2060           & 0.2726          & 0.1695          & 0.1952          \\ \cmidrule(l){2-8}
                             & SASRec      & 0.4145          & 0.2815          & 0.3148          & 0.4702          & 0.3042          & 0.3456          \\
                             & SSE-PT      & 0.3924	& 0.2595	& 0.2926   & 0.4931	& 0.3129	& 0.3578\\
                             & BERT4Rec    & 0.3650           & 0.2408          & 0.2718          & 0.4243          & 0.2718          & 0.3098          \\
                             & FMLP-Rec     & 0.4165          & \underline{0.2898}    & \underline{0.3215}    & \underline{0.5206}    & \underline{0.3476}    & \underline{0.3908}    \\ \cmidrule(l){2-8}
                             & PJFNN       & 0.2707          & 0.1478          & 0.1781          & 0.2636          & 0.1348          & 0.1666          \\
                             & BPJFNN      & 0.3791          & 0.2172          & 0.2573          & 0.4261          & 0.2273          & 0.2764          \\
                             & IPJF        & 0.3403          & 0.1743          & 0.2151          & 0.3776          & 0.1996          & 0.2435          \\
                             & PJFFF       & 0.3425          & 0.2126          & 0.2448          & 0.3715          & 0.1998          & 0.2422          \\
                             & DPGNN       & \underline{0.4521}    & 0.2696          & 0.3148          & 0.4409          & 0.2599          & 0.3047          \\ \cmidrule(l){2-8}
                             & ReSeq        & \textbf{0.7597*} & \textbf{0.4945*} & \textbf{0.5606*} & \textbf{0.7809*} & \textbf{0.5099*} & \textbf{0.5775*} \\ \bottomrule
\end{tabular}}}
\end{table*}

\begin{table*}[]
\caption{Performance comparison of different methods on two datasets in the questioner-answerer matching scenario. The best-performed and the second-best-performed methods are indicated in bold and underlined font, respectively. ``*'' denotes that the improvements are significant at the level of 0.01 with paired $t$-test.}
\label{tab:qa_comparison}
\resizebox{0.82\linewidth}{!}{
\setlength{\tabcolsep}{2.7mm}{
\begin{tabular}{@{}cccccccc@{}}
\toprule
\multirow{2}{*}{Dataset}       & Perspective & \multicolumn{3}{c}{Questioners}                     & \multicolumn{3}{c}{Answerers}                       \\ \cmidrule(l){2-2} \cmidrule(l){3-5} \cmidrule(l){6-8}
 & Metric   & HR@5   & MRR@5  & NDCG@5 & HR@5   & MRR@5  & NDCG@5 \\ \midrule
\multirow{8}{*}{StackOverflow} & BPR         & 0.3262          & 0.2344          & 0.2572          & 0.2653          & 0.1825          & 0.2029          \\
 & LFRR     & 0.3051 & 0.2113 & 0.2346 & 0.2391 & 0.1522 & 0.1737 \\
 & NeuMF    & 0.3218 & 0.2351 & 0.2567 & 0.2563 & 0.1781 & 0.1975 \\
 & LightGCN & 0.3352 & 0.2359 & 0.2606 & 0.2780  & 0.1881 & 0.2103 \\ \cmidrule(l){2-8}
 & SASRec   & 0.4336 & 0.3193 & 0.3478 & 0.2672 & 0.1780  & 0.2001 \\
 & SSE-PT   & 0.4064	& 0.2901	& 0.3190	& 0.2857	& 0.1881	& 0.2123  \\
 & BERT4Rec & 0.3621 & 0.2501 & 0.2779 & 0.2456 & 0.1548 & 0.1772 \\
                               & FMLP-Rec     & \underline{0.4594}    & \textbf{0.3428} & \textbf{0.3718} & \underline{0.2957}    & \underline{0.2026}    & \underline{0.2257}    \\ \cmidrule(l){2-8}
                               & ReSeq        & \textbf{0.4868*} & \underline{0.3235}    & \underline{0.3640}     & \textbf{0.4019*} & \textbf{0.2432*} & \textbf{0.2825*} \\ \midrule
\multirow{8}{*}{AskUbuntu}     & BPR         & 0.1566          & 0.1106          & 0.1220           & 0.1299          & 0.0914          & 0.1009          \\
 & LFRR     & 0.1590  & 0.1003 & 0.1149 & 0.0935 & 0.0622 & 0.0699 \\
 & NeuMF    & 0.2351 & 0.1592 & 0.1780  & 0.1323 & 0.0992 & 0.1074 \\
 & LightGCN & 0.1606 & 0.1128 & 0.1247 & 0.1245 & 0.0906 & 0.0990 \\ \cmidrule(l){2-8}
  & SASRec   & 0.2663 & 0.1698 & 0.1937 & 0.1588 & 0.1076 & 0.1202 \\
  & SSE-PT   & 0.2418	& 0.1584	& 0.1791	& 0.1540	& 0.0994	& 0.1129 \\
 & BERT4Rec & 0.2294 & 0.1404 & 0.1626 & 0.1143 & 0.0669 & 0.0786 \\
                               & FMLP-Rec     & \underline{0.2706}    & \underline{0.1818}    & \underline{0.2039}    & \underline{0.1627}    & \underline{0.1121}    & \underline{0.1246}    \\ \cmidrule(l){2-8}
                               & ReSeq        & \textbf{0.5259*} & \textbf{0.2848*} & \textbf{0.3447*} & \textbf{0.2254*} & \textbf{0.1455*} & \textbf{0.1653*} \\
                               \bottomrule
\end{tabular}}}
\end{table*}

\subsubsection{Baseline Models}
We compare ReSeq with the following baseline models (key configurations are in parentheses):
\begin{itemize}
\item {\textbf{BPR}~\cite{rendle2009bpr}} is a matrix factorization-based collaborative filtering method, which is optimized by Bayesian Personalized Ranking (BPR) loss (embedding size=64).
\item {\textbf{LFRR}~\cite{neve2019latent}} applies the latent factor model for reciprocal recommendation, which utilizes two latent factor models to model the dual perspectives of both parties (embedding size=64).
\item {\textbf{NeuMF}~\cite{he2017neural}} combines matrix factorization and neural networks to improve the effectiveness of personalized recommendation (MLP hidden sizes= [128,64], dropout ratio=0.1).
\item {\textbf{LightGCN}~\cite{he2020lightgcn}} is a collaborative filtering method based on lightweight graph convolutional networks, which simplifies graph convolution operations and achieves competitive performance on user-item recommendation (embedding size=64, number of layers=2).
\item {\textbf{SASRec}~\cite{kang2018self}} exploits a unidirectional Transformer-based neural network to model the item sequences and predict the next item (number of layers=2, number of attention heads=2, dropout ratio=0.5).
\item {\textbf{SSE-PT}~\cite{wu2020sse}} combines user and item embeddings as input of the Transformer model for personalized sequential recommendation (number of layers=2, number of attention heads=2, dropout ratio=0.5).
\item {\textbf{BERT4Rec}~\cite{sun2019bert4rec}} adopts a bidirectional Transformer model and combines it with a mask prediction task for the modeling of item sequences (number of layers=2, number of attention heads=2, fine-tuning ratio=0.5)
\item {\textbf{FMLP-Rec}~\cite{zhou2022filter}} proposes an all-MLP model with learnable filters, which ensures efficiency and reduces noisy signals (number of layers=2, dropout ratio=0.5).
\item {\textbf{PJFNN}~\cite{zhu2018person}} is a person-job fit model based on CNN to predict the matching degree between resumes and job descriptions (max sentence number=20, max sentence length=30).
\item {\textbf{BPJFNN}~\cite{qin2018enhancing}} leverages a bidirectional LSTM network to encode resumes and job descriptions for person-job fit prediction (max sentence length=512, number of layers=1, hidden size=32, dropout ratio=0.1)
\item {\textbf{IPJF}~\cite{le2019towards}} learns the matching propensities between candidates and recruiters based on various behavior labels in recruitment scenarios (max sentence number=20, max sentence length=30, BERT embedding size=768).
\item {\textbf{PJFFF}~\cite{jiang2020learning}} fuses the user explicit and implicit representations. The implicit representations are modeled by LSTM based on historical user records (hidden size=32, historical item length=100, BERT embedding size=768, BERT output size=32).
\item {\textbf{DPGNN}~\cite{yang2022modeling}} proposes a dual-perspective graph representation learning method that models recruitment as a two-way selection process (number of layers=2, BERT embedding size=768, BERT output size=32, contrastive loss weight=0.05, temperature=0.2).
\end{itemize}
Our baselines can be categorized into collaborative filtering models (\ie BPR, LFRR, NeuMF, LightGCN), sequential recommendation models (\ie SASRec, SSE-PT, BERT4Rec, FMLP-Rec), and person-job fit models (\ie PJFNN, BPJFNN, IPJF, PJFFF, DPGNN). Among these, the person-job fit models are exclusively applicable in online recruitment scenarios where users possess comprehensive textual information.

\subsubsection{Evaluation and Implementation Details}
Following ~\cite{yang2022modeling}, we employ three widely used metrics, Hit Ratio (HR@$k$), Normalized Discounted Cumulative Gain (NDCG@$k$), and Mean Reciprocal Rank (MRR@$k$) to evaluate ranking tasks from dual perspectives simultaneously, where $k$ is set to 5.
Concretely, for each positive instance interaction $(u_i, {S}^T_{u_i}, v_j, {S}^T_{v_j})$, we randomly sample $100$ users from $\mathcal{U}$ and $\mathcal{V}$ as dual-perspective negative instances, respectively. Then truncate their behavior sequences to the current interaction time for the evaluation to prevent data leakage~\cite{zhao2022revisiting}.

For all the compared methods, we use Adam for optimization and grid search for hyperparameter tuning to get optimal performance. The learning rates of all models are tuned in \{0.005, 0.001, 0.0005, 0.0003, 0.0001, 0.00001\}.
To ensure a fair comparison, the dimensions of embeddings are uniformly set to $64$. Besides, the weight decay coefficient is set to $1e-5$. And we select the best models by early stopping with the patience of $10$ epochs.
All models are implemented based on a popular open-source recommendation library \textsc{RecBole}~\cite{zhao2021recbole,zhao2022recbole}.

\begin{table*}[]
\caption{Ablation study of our method in two scenarios. The best-performed methods are indicated in bold.}
\label{tab:ablation_study}
\resizebox{0.85\linewidth}{!}{
\setlength{\tabcolsep}{2.5mm}{
\begin{tabular}{@{}clcccccc@{}}
\toprule
\multirow{2}{*}{Dataset} &
  Perspective &
  \multicolumn{3}{c}{Candidates/Questioners} &
  \multicolumn{3}{c}{Recruiters/Answerers} \\ \cmidrule(l){2-2} \cmidrule(l){3-5} \cmidrule(l){6-8}
 & 
  Metric &
  \multicolumn{1}{c}{HR@5} &
  \multicolumn{1}{c}{MRR@5} &
  \multicolumn{1}{c}{NDCG@5} &
  \multicolumn{1}{c}{HR@5} &
  \multicolumn{1}{c}{MRR@5} &
  \multicolumn{1}{c}{NDCG@5} \\ \midrule
\multirow{5}{*}{Design} &
  ReSeq &
  \textbf{0.4435} &
  \textbf{0.2554} &
  \textbf{0.3020} &
  \textbf{0.3722} &
  \textbf{0.1994} &
  \textbf{0.2420} \\
 & \ \ w/o DSE & 0.4224 & 0.2462 & 0.2898 & 0.3409 & 0.1806 & 0.2201 \\
 & \ \ w/o MASK & 0.4305 & 0.2510  & 0.2954 & 0.3406 & 0.1836 & 0.2223 \\
 & \ \ w/o TSA & 0.4139 & 0.2368 & 0.2806 & 0.3315 & 0.1798 & 0.2172 \\
 & \ \ w/o SD   & 0.4156 & 0.2372 & 0.2813 & 0.3446 & 0.1889 & 0.2273 \\ \midrule
\multirow{5}{*}{StackOverflow} &
  ReSeq &
  \textbf{0.4868} &
  \textbf{0.3235} &
  \textbf{0.3640} &
  \textbf{0.4019} &
  \textbf{0.2432} &
  \textbf{0.2825} \\
 & \ \ w/o DSE  & 0.4787 & 0.3102 & 0.3521 & 0.3964 & 0.2391 & 0.2780  \\
 & \ \ w/o MASK & 0.4764 & 0.3114 & 0.3524 & 0.3942 & 0.2387 & 0.2772 \\
 & \ \ w/o TSA & 0.4710  & 0.3128 & 0.3521 & 0.3697 & 0.2239 & 0.2600   \\
 & \ \ w/o SD   & 0.4787 & 0.3154 & 0.3560  & 0.3935 & 0.2377 & 0.2763 \\ \bottomrule
\end{tabular}}}
\end{table*}

\begin{figure*}[]
	\centering
	\includegraphics[width=0.495\linewidth]{./fig/lambda.pdf}
 \includegraphics[width=0.495\linewidth]{./fig/mu.pdf}
	\caption{Performance comparison w.r.t. micro teacher loss coefficient $\lambda$ and self-distillation loss coefficient $\mu$.}
	\label{fig:parameters}
\end{figure*}


\subsection{Experiment Results}
The results of different methods in the recruitment scenario are shown in Table~\ref{tab:recruit_comparison}, including baseline models dedicated to person-job fit (with extra text data). Table~\ref{tab:qa_comparison} shows the comparison between our model and baseline models in the questioner-answerer matching scenario. Based on these results, we can find:

For the four collaborative filtering methods, models based on GCN (\ie LightGCN) or MLP (\ie NeuMF) for hybrid interaction tend to exhibit superior performance in comparison to simple matrix factorization method (\ie BPR) and latent factor model (\ie LFRR). 
It indicates that increased granularity in interactions has a positive impact on the accuracy of recommendations, even in the case of static representations.
As for sequential recommendation methods, it can be observed that they generally achieve better performance compared to conventional collaborative filtering methods. We suppose this can be attributed to their ability to dynamically model user preferences based on the behavior sequences. 
Furthermore, FMLP-Rec attains considerable results compared to transformer-based models (\ie SASRec, SSE-PT, BERT4Rec). One possible reason for this is that the proposed learnable filtering layer has the ability to efficiently decrease noisy signals.

Additionally, in the recruitment scenario, there are also results of person-job fit methods. We can observe that models that rely exclusively on user resumes and job descriptions (\ie PJFNN and BPJFNN) are frequently subject to severe limitations caused by the quality and volume of textual data. 
These models perform better on Technology datasets that contain higher quality and larger volumes of text data compared to that Design and Sale datasets.
Among the hybrid models that combine both text and collaborative behavior signals (\ie IPJF, PJFFF, DPGNN), DPGNN achieves the best performance, which is inseparable from its modeling capacity of the two-way selection process in reciprocal recommendation.

Overall, compared with all baseline methods, it can be seen that our approach ReSeq achieves the best results in almost all the cases across five datasets from two scenarios. 
Different from these baselines, we formulate reciprocal recommendation as a distinctive sequence matching task, where matching prediction is performed based on the dynamic behavior sequences of both parties. 
And for lower time consumption, we introduce the self-distillation techniques to distill knowledge from the micro to macro level matching, which greatly speeds up the similarity calculation. As a result, our ReSeq is an effective and efficient reciprocal recommendation method.



\subsection{Ablation Study}
The major contribution of our method ReSeq lies in four parts: decomposed embedding, behavior sequence encoding, matching prediction, and efficiency improvement. 
In order to verify the effectiveness of specific designs for each component, we conduct an ablation study on Design and StackOverflow datasets to verify their contributions. 
Specifically, we consider the following four variants of ReSeq:
\begin{itemize}
\item {\underline{ReSeq w/o DSE}} removes the decomposition and sharing operations of embedding matrices (Eqn.~\eqref{eq:decomp}) and directly learns four embedding matrices for users.
\item {\underline{ReSeq w/o MASK}} 
replaces the different mask matrices for encoding user behavior sequences in Section~\ref{sec:seq_encode} with the same bidirectional mask matrix.
\item {\underline{ReSeq w/o TSA}} exchanges the time-sensitive attention aggregation operation in micro-level matching (Eqn.~\eqref{eq:agg}) into mean aggregation.
\item {\underline{ReSeq w/o SD}}  removes the self-distillation loss.
\end{itemize}

From the results in Table~\ref{tab:ablation_study}, it is evident that removing any of the above components would result in a decline in the overall effect. This indicates that all of the components in ReSeq play a vital role in delivering precise recommendations.


\subsection{Further Analysis}
\subsubsection{Hyperparameter Tuning}
In this section, we assess the robustness of our model and conduct a comprehensive analysis of two crucial loss function coefficients. The performance of the two datasets is presented in Figure~\ref{fig:parameters}. 
For each loss coefficient, we show the NDCG@5 of the two perspectives separately and calculate the average metrics of the two perspectives to reflect the overall trend.

For the micro teacher loss coefficient $\lambda$, we tune it within the range from 1 to 10 and additionally designate it as 0.5 to investigate the ramifications of inadequate teacher preparation. 
As shown in Figure~\ref{fig:parameters}, with the increasing of coefficient $\lambda$, the model's performance gradually improves. At the value of 5, ReSeq achieves the best performance, and further increment of $\lambda$ does not result in significant performance degradation, which fully reflects our method's robustness.

For the self-distillation loss coefficient $\mu$, we tune it in the ranges of \{0.0001, 0.001, 0.003, 0.005, 0.01, 0.05, 0.1\}. 
As shown in Figure~\ref{fig:parameters}, our model achieves the best results on the two datasets when $\mu$ is set to 0.005 and 0.01, respectively. A self-distillation loss coefficient that is too small cannot effectively transfer knowledge between micro and macro levels, while a too-large loss coefficient may cause the model to disregard ranking optimization.


\subsubsection{Deployment Efficiency Analysis}
To investigate the impact of micro-to-macro self-distillation on time latency in real-world deployment, we compare the efficiency of matching prediction between SASRec, ReSeq, and ReSeq w/o Self-Distillation on two datasets from different scenarios, \ie Design and StackOverflow.
All experiments are carried out on a Linux system machine with an Intel(R) Xeon(R) CPU E5-2630 v4 @ 2.20GHz CPU and a 12GB NVIDIA TITAN V GPU. For a fair comparison, we use the PyTorch framework to implement all methods and ensure that the common parameters such as batch size are uniform.
Specifically, we count the average matching prediction time of each batch of the three methods on the test sets. Besides, we maintain process synchronization through the \texttt{synchronize} function in PyTorch to obtain accurate time latency.

As shown in Figure~\ref{fig:efficiency}, our ReSeq has a latency slightly higher than SASRec. The additional consumption is an inevitable time increase brought by dual-perspective matching in reciprocal recommendation to achieve effective sequence matching.
However, without self-distillation, the fine-grained micro-level matching prediction would result in decades of increase in time consumption, which is unacceptable in actual deployment. In contrast, the micro-to-macro self-distillation enables our ReSeq to be an efficient and effective reciprocal recommendation method.

\begin{figure}[]
\centering
\includegraphics[width=1.0\linewidth]{./fig/efficiency-t.pdf}
\caption{Per-batch matching prediction latency.}
\label{fig:efficiency}
\end{figure}

\begin{figure}[]
\centering
\includegraphics[width=1.0\linewidth]{./fig/addtext.pdf}
\caption{Performance after adding text-based matching.}
\label{fig:addtext}
\end{figure}

\subsubsection{Text Enhancement}
Out of all, we conduct a text enhancement experiment on recruitment datasets. In detail, we directly incorporate a text matching score derived from pre-trained twin-tower BERT with an additional prediction layer into the macro-level matching prediction (Eqn.~\eqref{eq:ma_score}). The results are presented in Figure~\ref{fig:addtext}. Observably, after combining with simple text enhancement, our method basically obtains performance improvements in both perspectives on the three datasets.
Particularly, ReSeq benefits more on the Technology dataset.
The potential explanation for this phenomenon is that the text quality and quantity vary prominently across different datasets. Thus, in the Technology dataset, whose text quality and quantity substantially surpass the other two datasets, the effect of text enhancement is more significant. 
The above results illustrate that our method can be improved by incorporating some supplementary information.
Certainly, this simple text enhancement method may impose some limitations, and in the future, we will take into account more rational approaches to merge text or other information.


